% !TeX root = ../main.tex
\chapter{Discussion générale}
\label{chap:analyse}

\vspace{-12pt}
Le tableau~\ref{tab:recap-labo} synthétise l'ensemble des grandeurs mesurées et leurs incertitudes estimées.

\vspace{-6pt}
\begin{table}[H]
  \centering
  \scriptsize
  \setlength{\tabcolsep}{2.5pt}
  \renewcommand{\arraystretch}{0.95}
  \begin{tabularx}{\textwidth}{l >{\raggedright\arraybackslash}X c c c}
    \toprule
    \textbf{Dispositif} & \textbf{Grandeur} & \textbf{Résultat mesuré} & \textbf{Référence / Condition} & \textbf{Incertitude / Écart} \\
    \midrule
    \multirow{2}{*}{CTN SEMI833} & Résistance ambiante $R_{\text{amb}}$ & \qty{97.20}{\kilo\ohm} & \qty{103.60}{\kilo\ohm} (\qty{20}{\degreeCelsius}) & Pente constructeur \\
     & Température ambiante $T_{\text{amb, NTC}}$ & \qty{21.54 \pm 0.80}{\degreeCelsius} & -- & Tolérance $\pm \qty{0.8}{\degreeCelsius}$ \\
    \midrule
    \multirow{5}{*}{Pt100 KP (chemisée)} & Résistance 4W ambiante & \qty{108.530}{\ohm} & \qty{100.00}{\ohm} (\qty{0}{\degreeCelsius}) & Multimètre 6½ digits \\
     & Température ambiante $T_{\text{amb, Pt100}}$ & \qty{22.16 \pm 0.15}{\degreeCelsius} & -- & Tolérance classe A ($\pm \qty{0.15}{\degreeCelsius}$) \\
     & Température LabVIEW & \qty{22.08 \pm 0.15}{\degreeCelsius} & \qty{22.16}{\degreeCelsius} & Écart : \qty{-0.08}{\degreeCelsius} \\
     & Résistance étalonnée $R_0$ (\qty{0}{\degreeCelsius}) & \qty{100.10 \pm 0.03}{\ohm} & \qty{100.00}{\ohm} & Écart : \qty{+0.10}{\ohm} (\qty{+0.1}{\percent}) \\
     & Biais de ligne 2W ($2 R_L$) & \qty{0.09}{\ohm} & \qty{0.00}{\ohm} (4W) & $\Delta T = \qty{+0.23}{\degreeCelsius}$ \\
    \midrule
    \multirow{2}{*}{Dynamique Pt100} & Constante de temps $\tau_{\text{fit}}$ & \qty{5.80 \pm 0.05}{\second} & \qty{5.80}{\second} (tangente) & $\text{RMSE} = \qty{0.09}{\degreeCelsius}$ \\
     & Temps de réponse à \qty{95}{\percent} & \qty{17.37 \pm 0.15}{\second} & $3\tau = \qty{17.40}{\second}$ & Écart : \qty{-0.2}{\percent} \\
    \midrule
    \multirow{2}{*}{Flux d'air CTN} & Palier standard ($N=497$) & \qty{41.13 \pm 0.13}{\degreeCelsius} & -- & Bruit : $\sigma = \qty{2.98}{\degreeCelsius}$ \\
     & Palier suréchantillonné ($N=352$) & \qty{42.83 \pm 0.17}{\degreeCelsius} & -- & Bruit : $\sigma = \qty{3.15}{\degreeCelsius}$ \\
    \midrule
    Auto-échauffement & Puissance $P$ / $\Delta T$ & \qty{0.061}{\milli\watt} & $\Delta T \approx \qty{0.09}{\degreeCelsius}$ & Négligeable ($< \qty{0.8}{\degreeCelsius}$) \\
    \bottomrule
  \end{tabularx}
  \caption{Synthèse des résultats expérimentaux et incertitudes associées.}
  \label{tab:recap-labo}
\end{table}

\vspace{-10pt}
\section{Compatibilité entre la mesure CTN et la mesure Pt100}
\label{sec:compatibilite}

À l'ambiance, la CTN donne $T_{\text{amb, NTC}} = \qty{21.54 \pm 0.80}{\degreeCelsius}$ et la Pt100 donne $T_{\text{amb, Pt100}} = \qty{22.16 \pm 0.15}{\degreeCelsius}$ ($\Delta T = \qty{0.62}{\degreeCelsius}$). L'incertitude combinée vaut $u_c = \sqrt{0.80^2 + 0.15^2} \approx \qty{0.81}{\degreeCelsius}$, soit un critère de compatibilité normalisé :
\begin{equation}
    C = \frac{\Delta T}{u_c} = \frac{0.62}{0.814} \approx \mathbf{0.76} \le 2
\end{equation}
Comme $C \le 2$, l'écart de \qty{0.62}{\degreeCelsius} est statistiquement justifié par les incertitudes : les deux capteurs sont en accord parfait.

\vspace{-4pt}
\section{Analyse des principales sources d'erreurs}
\label{sec:sources_erreurs}

\begin{enumerate}[leftmargin=1.0cm,topsep=1pt,itemsep=0.5pt,parsep=0pt]
    \item \textbf{Résistance de ligne 2 fils sur la Pt100 ($+\qty{0.23}{\degreeCelsius}$) :} Le câble de \qty{1}{\meter} ajoute $2 R_{\text{ligne}} = \qty{0.09}{\ohm}$. Avec $S = \qty{0.385}{\ohm\per\degreeCelsius}$, cela provoque un biais systématique de $+\qty{0.23}{\degreeCelsius}$ (supérieur à la tolérance classe A de $\pm \qty{0.15}{\degreeCelsius}$), démontrant l'obligation du câblage 4 fils.
    \item \textbf{Auto-échauffement CTN ($\approx \qty{0.09}{\degreeCelsius}$) :} Sous \qty{14}{\volt} avec $R = \qty{462.35}{\kilo\ohm}$, la puissance dissipée ($P = \qty{0.061}{\milli\watt}$) limite l'échauffement interne à $\Delta T \approx \qty{0.09}{\degreeCelsius}$ (\qty{11}{\percent} de la tolérance), validant le dimensionnement du pont.
    \item \textbf{Turbulence du flux d'air et tenue manuelle :} Le suréchantillonnage élimine le bruit de numérisation, mais l'écart-type résiduel ($\sigma \approx \qty{3}{\degreeCelsius}$) est dominé par les tourbillons réels de l'air chaud. Le décalage moyen ($\approx \qty{1.7}{\degreeCelsius}$) s'explique par la tenue manuelle non guidée du sèche-cheveux et de la perle.
\end{enumerate}
